\documentclass[11pt,a4paper]{article}

\usepackage[french]{babel}
\usepackage[utf8]{inputenc}
\usepackage[T1]{fontenc}
\usepackage{lmodern}
\usepackage[margin=2.5cm]{geometry}
\usepackage{amsmath}
\usepackage{booktabs}
\usepackage{longtable}
\usepackage{array}
\usepackage{xcolor}
\usepackage{hyperref}
\usepackage{fancyhdr}

\definecolor{solidateal}{RGB}{15,76,79}
\definecolor{gris}{RGB}{90,90,90}

\hypersetup{
  colorlinks=true,
  linkcolor=solidateal,
  urlcolor=solidateal,
  pdftitle={SOLIDA -- Registre des formules de calcul},
  pdfauthor={SOLIDA}
}

\pagestyle{fancy}
\fancyhf{}
\fancyhead[L]{\small\textcolor{gris}{SOLIDA -- Registre des formules de calcul}}
\fancyhead[R]{\small\textcolor{gris}{\thepage}}
\renewcommand{\headrulewidth}{0.4pt}

% Chemins de fichiers : coupables apres "/" et "::" pour eviter les debordements.
\newcommand{\brk}{\allowbreak}
\newcommand{\ecran}[1]{\noindent\textbf{Écran(s) concerné(s) :} #1\par}
\newcommand{\fichier}[1]{\noindent\textbf{Implémentation :} \tiny\ttfamily#1\normalfont\normalsize\par}
\newenvironment{formule}{\begin{quote}}{\end{quote}}

% Colonnes de tableau qui retournent a la ligne au lieu de deborder.
\newcolumntype{L}[1]{>{\raggedright\arraybackslash}p{#1}}

\title{\textcolor{solidateal}{\Huge SOLIDA} \\[0.3em] \Large Registre des formules de calcul}
\author{}
\date{Dernière mise à jour : \today}

\begin{document}

\maketitle

\section*{Objet de ce document}

Ce document répertorie \textbf{toutes} les formules de calcul utilisées par SOLIDA pour produire
un score, une décision, un plafond de crédit ou une échéance -- sans exception. Pour chacune, il
précise : l'écran où le résultat apparaît, chaque variable et chaque constante en jeu avec sa
valeur actuelle, la formule telle qu'elle est implémentée aujourd'hui, pourquoi elle a été
choisie ainsi (y compris quand la réponse honnête est « aucune source, choix arbitraire »), et ce
que le modèle de scoring réel changera -- ou ne changera pas -- une fois entraîné.

C'est la référence unique pour comprendre ou modifier un calcul sans relire le code source. Le
code lui-même ne porte que des renvois courts vers ce document : le détail et la justification
vivent ici, pas dans des commentaires longs qui alourdissent la lecture du code.

\textbf{Principe transversal, valable pour toutes les formules ci-dessous sauf mention contraire :}
le modèle de scoring réel, une fois entraîné, améliore uniquement la probabilité de défaut $p$
qu'il produit. Il ne calibre jamais tout seul les paramètres de politique de risque (marges,
seuils, coefficients de modulation) qui transforment ensuite cette probabilité en score, en
décision ou en plafond -- ce sont des arbitrages de la coopérative, qui se calibrent par
rétro-test sur des cohortes réelles de remboursement, pas par l'entraînement du modèle. Chaque
section précise l'exception quand elle existe.

\textbf{Origine des valeurs actuelles.} Une partie des constantes ci-dessous (grille de décision,
crédit progressif, conditions de réexamen) proviennent d'un prototype de calcul,
\texttt{simulateur/\brk decision.py}, écrit pour démontrer le mécanisme avant que le vrai backend
n'existe, puis retiré du dépôt une fois son rôle de démonstration rempli (trace complète dans
\texttt{03-MODELE/\brk 09-lecons-prototype-simulateur.md}). Ce prototype n'a jamais prétendu être
calibré sur des données réelles : c'est un point de départ illustratif, pas une vérité mesurée.
Chaque section indique explicitement quand une valeur a cette origine, et quand elle n'a
\textbf{aucune} source documentée du tout.

\tableofcontents

\newpage

% ============================================================
\section{Transformation probabilité → score (échelle PDO)}
\label{sec:scorecard}

\ecran{Résultat de scoring (jauge et nombre de points), fiche de justification.}
\fichier{backend/\brk solida/\brk domain/\brk rules/\brk scorecard.py::\brk{}calculer\_score}

\subsection*{Variables}

\begin{tabular}{@{}lL{10cm}@{}}
\toprule
Symbole & Signification \\
\midrule
$p$ & Probabilité de défaut produite par le modèle de scoring \\
\bottomrule
\end{tabular}

\subsection*{Constantes}

\begin{tabular}{@{}L{3.2cm}rL{7.4cm}@{}}
\toprule
Constante & Valeur & Pourquoi cette valeur \\
\midrule
$\mathit{PDO}$ (points pour doubler les cotes) & $20$ &
  Reprend l'exemple canonique de la littérature scorecard (popularisé par les manuels de
  \emph{credit scoring} et repris comme valeur par défaut dans la plupart des outils de
  scorecard), pas une calibration propre à SOLIDA. \\
$\mathit{score\_ref}$ & $600$ & Idem : point d'ancrage canonique de la même littérature. \\
$\mathit{odds\_ref}$ & $50$ & Idem : rapport de cotes canonique associé à $\mathit{score\_ref}=600$
  dans le même exemple de référence. \\
$\mathit{score\_min}$ & $300$ & Convention du score FICO américain (300--850), reprise pour sa
  familiarité, sans lien avec la distribution réelle du futur modèle. \\
$\mathit{score\_max}$ & $850$ & Idem. \\
\bottomrule
\end{tabular}

\subsection*{Formule}
\begin{formule}
\begin{align}
\mathit{Facteur} &= \frac{\mathit{PDO}}{\ln 2} \\
\mathit{Décalage} &= \mathit{score\_ref} - \mathit{Facteur} \times \ln(\mathit{odds\_ref}) \\
\mathit{score} &= \operatorname{clip}\left(\mathit{Décalage} + \mathit{Facteur} \times \ln\!\left(\frac{1-p}{p}\right),\ \mathit{score\_min},\ \mathit{score\_max}\right)
\end{align}
\end{formule}

Convention bancaire : score élevé $=$ risque faible, d'où le rapport $(1-p)/p$ et non son inverse.

\subsection*{Pourquoi cette formule}
C'est la transformation PDO standard de la scorecard bancaire (logistique $\to$ score linéaire),
qui garantit qu'une variation constante du score correspond toujours au même facteur
multiplicatif sur le rapport de cotes, quel que soit le niveau de risque de départ. C'est ce qui
rend le score lisible et comparable d'un dossier à l'autre.

\subsection*{Statut et rôle du modèle réel}
\textbf{Exception au principe transversal} : contrairement à la grille et au crédit progressif,
$\mathit{PDO}$, $\mathit{score\_ref}$ et $\mathit{odds\_ref}$ \emph{doivent} être recalibrés dès
que le modèle réel existe -- l'échelle n'a de sens que rapportée à la distribution réelle de $p$
que ce modèle produit. Un score mal calibré resterait affiché de façon trompeuse même si la
décision sous-jacente (section~\ref{sec:grille}) reste correcte, puisqu'elle raisonne directement
sur $p$, jamais sur le score affiché. $\mathit{score\_min}$ et $\mathit{score\_max}$ n'ont pas
besoin de changer : ce sont de simples bornes d'affichage.

\newpage

% ============================================================
\section{Invariant de cohérence de la décomposition}
\label{sec:invariant}

\ecran{Aucun -- vérification interne avant l'enregistrement de toute décision.}
\fichier{backend/\brk solida/\brk domain/\brk rules/\brk scorecard.py::\brk{}verifier\_invariant\_decomposition}

\subsection*{Constante}

\begin{tabular}{@{}L{3.2cm}rL{7.4cm}@{}}
\toprule
Constante & Valeur & Pourquoi cette valeur \\
\midrule
\texttt{TOLERANCE\_\brk{}INVARIANT} & $1{,}0$ point & Marge technique pour l'arrondi
  flottant accumulé en sommant plusieurs contributions, pas un seuil métier : assez large pour
  ne jamais déclencher sur un arrondi légitime, assez serrée pour rejeter toute vraie
  incohérence de calcul. \\
\bottomrule
\end{tabular}

\subsection*{Formule}
\begin{formule}
\[
\left|\ \mathit{points\_de\_base} + \sum_j \mathit{points}_j - \mathit{score}\ \right| < \mathit{TOLERANCE\_INVARIANT}
\]
\end{formule}
Si l'inégalité est violée, le scoring est \textbf{rejeté} (exception \texttt{InvariantScoreViole})
plutôt que rendu de façon incohérente.

\subsection*{Pourquoi ce mécanisme}
Une fiche de justification dont les points ne somment pas exactement au score détruirait la
crédibilité du produit devant un sociétaire ou un auditeur. Mieux vaut refuser un scoring que
l'afficher faux.

\subsection*{Statut et rôle du modèle réel}
Aucun changement attendu : cette vérification reste valable quelle que soit la source des
contributions, apprises ou factices.

\newpage

% ============================================================
\section{Décomposition en points (« Facteurs déterminants »)}
\label{sec:decomposition}

\ecran{Résultat de scoring (graphique « Facteurs déterminants »), fiche de justification.}
\fichier{backend/\brk solida/\brk domain/\brk rules/\brk scorecard.py::\brk{}decomposer\_en\_points}

\subsection*{Variables}

\begin{tabular}{@{}lL{10cm}@{}}
\toprule
Symbole & Signification \\
\midrule
$\beta_0$ & Constante (log-cotes de base, hors effet des variables) \\
$f_j$ & Contribution en log-cotes de la variable $j$ \\
$\mathit{Facteur}, \mathit{Décalage}$ & Identiques à la section~\ref{sec:scorecard} \\
\bottomrule
\end{tabular}

\subsection*{Formule}
\begin{formule}
\begin{align}
\mathit{points\_de\_base} &= \mathit{Décalage} + \mathit{Facteur} \times \beta_0 \\
\mathit{points}_j &= \mathit{Facteur} \times f_j
\end{align}
\end{formule}

\subsection*{Pourquoi cette formule}
La transformation probabilité $\to$ score étant affine, chaque terme du log-cotes se met à
l'échelle indépendamment par le même facteur : la décomposition en points hérite donc exactement
des mêmes propriétés que le score global, ce qui rend l'invariant de la section~\ref{sec:invariant}
possible.

\subsection*{Statut et rôle du modèle réel}
Le modèle réel doit fournir des $f_j$ qui soient de véritables contributions en log-cotes
(nativement additives pour un modèle de type EBM/GAM, ou approximées par des valeurs de Shapley
pour un autre type de modèle). Voir section~\ref{sec:contributions-factices} pour l'état actuel de
$\beta_0$ et des $f_j$, aujourd'hui entièrement fictifs.

\newpage

% ============================================================
\section{Contributions factices du modèle de substitution}
\label{sec:contributions-factices}

\ecran{Résultat de scoring (graphique « Facteurs déterminants »), fiche de justification.}
\fichier{backend/\brk solida/\brk adapters/\brk ml/\brk modele\_constant.py::\brk{}ModeleConstant.contributions ;
backend/\brk solida/\brk application/\brk use\_cases/\brk scorer\_demande.py (calcul de beta\_0)}

\subsection*{Contexte}
\texttt{ModeleConstant} (le modèle de substitution utilisé tant qu'aucun modèle réel n'est
entraîné) renvoie une probabilité \textbf{fixe} $p = 0{,}09$ -- le taux de créances en souffrance
cible du générateur de données, pas une valeur inventée. Comme $p$ ne varie jamais avec le
dossier, $\beta_0$ est recalculé à chaque scoring pour absorber exactement l'écart :
\begin{formule}
\[
\beta_0 = \ln\!\left(\frac{1-p}{p}\right) - \sum_j f_j
\]
\end{formule}
Conséquence mathématique importante : \textbf{le score et la décision ne dépendent jamais} de ce
que \texttt{contributions()} renvoie, quelle que soit sa valeur. C'est ce qui rend sûr d'y mettre
une décomposition non apprise sans risquer de fausser une décision réelle.

\subsection*{Constantes -- coefficients choisis à la main, pas appris}

\begin{tabular}{@{}L{5.2cm}rr@{}}
\toprule
Variable & Coefficient & Point de référence \\
\midrule
Régularité des dépôts (12 mois) & $0{,}12$ & 6 mois \\
Taux d'endettement & $-1{,}4$ & $0{,}35$ \\
Nombre d'incidents antérieurs & $-0{,}5$ & $0$ \\
Ancienneté sociétaire (mois) & $0{,}01$ & 36 mois \\
Épargne rapportée au revenu & $0{,}6$ & $0{,}3$ \\
Épargne rapportée au montant demandé & $0{,}5$ & $0{,}5$ \\
\bottomrule
\end{tabular}

\subsection*{Pourquoi ces valeurs précises}
\textbf{Aucune source, aucune calibration.} Seul le \emph{signe} de chaque coefficient reprend
l'intuition métier attendue (même intuition que \texttt{simulateur/\brk config/\brk config.yaml},
jamais ses valeurs). Les points de référence sont choisis pour être des valeurs médianes
plausibles d'un dossier moyen, à l'œil, sans mesure. C'est une démonstration délibérément
assumée comme non apprise, pas une estimation.

\subsection*{Formule (par variable)}
\begin{formule}
\[
f_j = \mathit{coefficient}_j \times (\mathit{valeur}_j - \mathit{référence}_j)
\]
\end{formule}

\subsection*{Pourquoi cette fonction existe}
Sans elle, l'écran « Facteurs déterminants » afficherait une liste vide (« Aucune contribution
significative identifiée »), ce qui est correct mais peu utile pour juger l'interface avant qu'un
modèle réel existe.

\subsection*{Statut et rôle du modèle réel}
\textbf{C'est la toute première chose à supprimer} dès qu'un modèle réel (EBM entraîné) existe.
\texttt{decomposer\_en\_points} (section~\ref{sec:decomposition}) doit alors recevoir les vraies
contributions log-cotes apprises, à la place de cette liste écrite à la main. Rien d'autre en
aval (mapper HTTP, graphique, fiche PDF) n'a besoin de changer : ils affichent déjà correctement
n'importe quelle décomposition non vide.

\newpage

% ============================================================
\section{Grille de décision (tranche)}
\label{sec:grille}

\ecran{Résultat de scoring (badge de décision, zones colorées de la jauge), écran de paramétrage
de la grille.}
\fichier{backend/\brk solida/\brk domain/\brk rules/\brk grille.py::\brk{}decider ;
frontend/\brk lib/\brk scorecard.ts::\brk{} trancheDepuisProbabilite (miroir, même formule)}

\subsection*{Variables}

\begin{tabular}{@{}lL{9.5cm}@{}}
\toprule
Symbole & Signification \\
\midrule
$p$ & Probabilité de défaut produite par le modèle \\
\bottomrule
\end{tabular}

\subsection*{Constantes}

\begin{tabular}{@{}L{3.2cm}rL{6.9cm}@{}}
\toprule
Constante & Valeur & Pourquoi cette valeur \\
\midrule
$\mathit{marge}$ & $0{,}15$ & Hérite du prototype retiré \texttt{decision.py} -- un choix
  illustratif pour démontrer le mécanisme, sans dérivation actuarielle sur des données
  réelles. \\
$\mathit{lgd}$ (perte en cas de défaut) & $0{,}75$ & Même origine que $\mathit{marge}$. \\
$m_a$ (multiplicateur accord) & $0{,}6$ & Même origine. \\
$m_v$ (multiplicateur vigilance) & $1{,}0$ & Même origine (le seuil économique lui-même, sans
  marge de tolérance). \\
$m_e$ (multiplicateur examen) & $1{,}6$ & Même origine. \\
$\mathit{plafond\_institutionnel}$ & $100\,000\,000$ FCFA & Maximum institutionnel unique,
  sans plafond par produit -- décision terrain actée
  (\texttt{docs/\brk continuite/\brk 2026-09-13-constat-claude-plafonds.md}). Champ de la
  grille depuis le 2026-09-14 (auparavant une constante Python en dur). \\
\bottomrule
\end{tabular}

\subsection*{Formule}
\begin{formule}
\begin{align}
\mathit{seuil} &= \frac{\mathit{marge}}{\mathit{marge} + \mathit{lgd}} \\[0.3em]
\mathit{tranche} &=
\begin{cases}
\text{ACCORD} & \text{si } p < \mathit{seuil} \times m_a \\
\text{ACCORD\_SOUS\_CONDITION} & \text{si } p < \mathit{seuil} \times m_v \\
\text{COMITÉ\_DE\_CRÉDIT} & \text{si } p < \mathit{seuil} \times m_e \\
\text{REFUS} & \text{sinon}
\end{cases}
\end{align}
\end{formule}

\subsection*{Pourquoi cette formule}
$\mathit{seuil}$ est le seuil économique classique qui égalise le coût espéré des deux erreurs
possibles : refuser un bon dossier (perte de marge) et accorder un mauvais dossier (perte en cas
de défaut). Les multiplicateurs créent des zones de tolérance autour de ce seuil plutôt qu'une
coupure binaire, pour laisser une marge de jugement humain (accord sous condition, comité de
crédit) avant le refus pur. La jauge du résultat de scoring recalcule ces mêmes seuils en score
(via la formule de la section~\ref{sec:scorecard}) pour que les zones colorées affichées
correspondent exactement aux frontières de décision -- aucune constante supplémentaire n'y est
introduite.

\textbf{Ce n'est pas une formule maison} : c'est exactement le seuil de décision bayésien optimal
sous coûts asymétriques établi par Elkan (2001), équation~2 -- $p^{*} = (c_{10}-c_{00}) /
(c_{10}-c_{00}+c_{01}-c_{11})$, qui se réduit à $c_{10}/(c_{10}+c_{01})$ quand les prédictions
correctes ne coûtent rien (hypothèse standard). $\mathit{marge}$ correspond à $c_{10}$ (coût
d'un bon dossier refusé, le manque à gagner) et $\mathit{lgd}$ à $c_{01}$ (coût d'un mauvais
dossier accepté, la perte en cas de défaut). Seules les \emph{valeurs} de $\mathit{marge}$ et
$\mathit{lgd}$ restent à calibrer sur des coûts réels (recherche du 2026-09-14, source
officielle) ; la formule elle-même n'a pas besoin d'être revue.

\textit{Source : C. Elkan, ``The Foundations of Cost-Sensitive Learning'', IJCAI-01, 2001,
\S1.3, équation (2) --
\url{https://cseweb.ucsd.edu/~elkan/rescale.pdf}.}

\subsection*{Statut et rôle du modèle réel}
Le modèle réel améliore uniquement $p$, l'entrée de \texttt{decider}. Il ne calibre jamais
$\mathit{marge}$, $\mathit{lgd}$ ni les multiplicateurs : ce sont des paramètres de politique de
risque de la coopérative, ajustables depuis l'écran de paramétrage de la grille, à rétro-tester
sur des cohortes réelles une fois qu'elles existeront.

\newpage

% ============================================================
\section{Plafond de crédit progressif}
\label{sec:progressif}

\ecran{Résultat de scoring (« Montant recommandé »).}
\fichier{backend/\brk solida/\brk domain/\brk rules/\brk progressif.py::\brk{}calculer\_plafond, \_modulation\_risque}

\subsection*{Variables}

\begin{tabular}{@{}lL{9.5cm}@{}}
\toprule
Symbole & Signification \\
\midrule
$p$ & Probabilité de défaut produite par le modèle \\
$M$ & Montant maximum déjà remboursé par le sociétaire (historique), ou \emph{néant} pour un
  primo-emprunteur \\
$D$ & Montant demandé pour ce nouveau crédit \\
\bottomrule
\end{tabular}

\subsection*{Constantes}

\begin{tabular}{@{}L{3.2cm}rL{6.5cm}@{}}
\toprule
Constante & Valeur & Pourquoi cette valeur \\
\midrule
$c$ (coefficient de progression) & $1{,}5$ & Hérite du prototype retiré \texttt{decision.py}, sans
  dérivation actuarielle sur des données réelles. \\
$\mathit{plancher}$ & 50\,000 FCFA & Même origine. \\
$\mathit{plafond\_produit}$ & ex. 3\,000\,000 FCFA & Résolu par produit de crédit (voir
  section~\ref{sec:echeance}), même origine, ajustable par la supervision. \\
$\mathit{plafond\_primo}$ & 150\,000 FCFA & Même origine. \\
$a$ (base de modulation) & $1{,}3$ & Même origine. \\
$b$ (pente de modulation) & $2{,}0$ & Même origine. \\
$m_{\min}$ & $0{,}4$ & Même origine. \\
$m_{\max}$ & $1{,}2$ & Même origine. \\
\bottomrule
\end{tabular}

\subsection*{Formule}

\textbf{Primo-emprunteur} ($M$ absent) :
\begin{formule}
\[
\mathit{montant\_recommandé} = \max\Big(\min(\mathit{plafond\_primo},\ D),\ \mathit{plancher}\Big)
\]
\end{formule}

\textbf{Sociétaire avec historique} :
\begin{formule}
\begin{align}
\mathit{modulation} &= \operatorname{clip}(a - b \times p,\ m_{\min},\ m_{\max}) \\
\mathit{base} &= \max(M \times c,\ \mathit{plancher}) \\
\mathit{montant\_recommandé} &= \max\Big(\min(\mathit{base} \times \mathit{modulation},\ \mathit{plafond\_produit},\ D),\ \mathit{plancher}\Big)
\end{align}
\end{formule}

\subsection*{Pourquoi cette formule}
Le crédit progressif (récompenser un historique de remboursement sain par un plafond plus élevé)
est un mécanisme d'incitation classique en microfinance, documenté depuis le modèle Grameen
(prêt progressif / \emph{step lending}). La modulation par $p$ ajoute une dimension que ce
mécanisme historique n'a pas seul : elle évite d'accorder une forte augmentation à un sociétaire
dont le profil de risque s'est dégradé entre deux cycles, même si son historique de remboursement
reste bon sur le papier -- cohérent avec la pratique de gestion de ligne de crédit bancaire, où le
score comportemental borne toujours la décision dans une règle explicite, jamais l'inverse.

\subsection*{Statut et rôle du modèle réel}
Le modèle réel améliore uniquement $p$, l'entrée de $\mathit{modulation}$. Il ne calibre jamais
$c$, $\mathit{plancher}$, $\mathit{plafond\_primo}$ ni $(a, b, m_{\min}, m_{\max})$ : ce sont des
paramètres de politique de risque du comité de crédit, à rétro-tester sur des cohortes réelles de
remboursement une fois qu'elles existeront en production -- un travail distinct de l'entraînement
du modèle de score.

\newpage

% ============================================================
\section{Palier accessible au prochain cycle}
\label{sec:trajectoire}

\ecran{Résultat de scoring (« Palier suivant accessible »).}
\fichier{backend/\brk solida/\brk domain/\brk rules/\brk progressif.py::\brk{}calculer\_trajectoire}

\subsection*{Constante}

\begin{tabular}{@{}L{3.2cm}rL{6.9cm}@{}}
\toprule
Constante & Valeur & Pourquoi cette valeur \\
\midrule
$\mathit{nb\_cycles}$ & $1$ & Voir justification détaillée ci-dessous : aucune institution
  sérieuse ne projette plusieurs cycles à l'avance. \\
\bottomrule
\end{tabular}

\subsection*{Formule}
\begin{formule}
\[
\mathit{palier} = \min\big(\mathit{montant\_recommandé} \times c \times \mathit{modulation},\ \mathit{plafond\_produit}\big)
\]
\end{formule}
où $c$ et $\mathit{modulation}$ sont identiques à la section~\ref{sec:progressif},
$\mathit{modulation}$ étant gelée sur la probabilité $p$ connue \emph{aujourd'hui}.

\subsection*{Pourquoi une projection sur un seul cycle}
Aucune institution de microcrédit sérieuse ne pré-calcule ni ne garantit un plafond à plusieurs
cycles d'écart. Les \emph{Guidelines for Establishing and Operating Grameen-Style Microcredit
Programs} (Grameen Foundation) l'interdisent même explicitement : chaque proposition de prêt est
évaluée et approuvée \emph{au moment du renouvellement}, sur la capacité de remboursement et
l'historique du cycle qui vient de s'achever -- jamais promise à l'avance -- et le code de
conduite du gestionnaire de centre y est explicite : \emph{« Does not give any false hope nor
lies to members »}, \emph{« Strictly keeps any word given to a loanee. »} La gestion de ligne de
crédit bancaire suit le même principe (score comportemental réévalué à chaque cycle réel, jamais
projeté). Chaque cycle supplémentaire empilerait une hypothèse de bon remboursement sur une
hypothèse de risque inchangé, de moins en moins défendable.

\subsection*{Statut et rôle du modèle réel}
Identique à la section~\ref{sec:progressif} : seul $p$ change avec le modèle réel. La fonction
reste paramétrable (\texttt{nb\_cycles}) si une projection plus longue est un jour explicitement
demandée par la coopérative à des fins d'incitation -- ce serait alors un choix produit assumé,
pas une pratique standard.

\newpage

% ============================================================
\section{Conditions de réexamen}
\label{sec:reexamen}

\ecran{Résultat de scoring (« Conditions de réexamen »), affiché pour toute tranche autre que
ACCORD.}
\fichier{backend/\brk solida/\brk domain/\brk rules/\brk progressif\_reexamen.py::\brk{}lister\_conditions\_reexamen}

\subsection*{Constantes -- cibles}

\begin{tabular}{@{}L{5.5cm}rL{4.6cm}@{}}
\toprule
Constante & Valeur & Pourquoi cette valeur \\
\midrule
$\mathit{régularité\_cible}$ & $75\,\%$ & Hérite du prototype retiré \texttt{decision.py}
  (voir \texttt{03-MODELE/\brk 09-lecons-\brk prototype-\brk simulateur.md}), sans calibration sur
  des cohortes réelles. \\
$\mathit{mois\_observation}$ & $3$ & Même origine. \\
$\mathit{ratio\_garantie\_cible}$ & $50\,\%$ & Même origine. \\
$\mathit{endettement\_seuil}$ & $33\,\%$ & \textbf{Plus une constante propre à cette règle} depuis
  le 2026-09-14 : reçu de \texttt{ParametresGrille.ratio\_endettement\_maximal} (section
  \ref{sec:capacite}), sourcé étude CEF-MF Lomé 2025 -- une seule valeur d'endettement dans tout
  le système. \\
\bottomrule
\end{tabular}

\subsection*{Formule (exemple : rattrapage de régularité)}
\begin{formule}
\[
\mathit{mois\_à\_rattraper} = \operatorname{round}\big((\mathit{régularité\_cible} - \mathit{régularité\_actuelle}) \times 12\big)
\]
\end{formule}
Les autres conditions (garantie, endettement, tendance d'épargne, caution déjà appelée) sont des
comparaisons directes au seuil, sans formule dédiée, traduites en phrase actionnable.

\subsection*{Pourquoi ce mécanisme}
Un refus ne doit jamais être une porte fermée : chaque condition est un levier concret et
vérifiable que le sociétaire peut activer avant le prochain cycle, plutôt qu'un « non » sans
recours.

\subsection*{Statut et rôle du modèle réel}
Indépendant du modèle de scoring : ces cibles sont des seuils métier, pas des sorties de modèle,
et le modèle n'influence pas quelles conditions sont listées -- seule la tranche décidée
(section~\ref{sec:grille}) détermine si ce bloc s'affiche.

\newpage

% ============================================================
\section{Cascade : choix du mode de calcul (socle vs enrichi)}
\label{sec:cascade}

\ecran{Résultat de scoring (mention « Score calculé avec l'historique du groupe de caution » ou
« profil individuel »).}
\fichier{backend/\brk solida/\brk domain/\brk rules/\brk cascade.py::\brk{}determiner\_mode}

\subsection*{Constantes}

\begin{tabular}{@{}L{5.5cm}rL{4.6cm}@{}}
\toprule
Constante & Valeur & Pourquoi cette valeur \\
\midrule
$\mathit{taille\_groupe\_minimale}$ & $5$ & Réponse terrain (\texttt{précision.txt}, question
  22 : « à partir de combien de membres peut-on réellement juger le comportement d'un groupe ? »,
  réponse « de 5 à 10 »). Relevé de $3$ à $5$ le 2026-09-14 ; aligné sur
  \texttt{SEUIL\_TAILLE\_GROUPE\_JUGEABLE} du catalogue d'entraînement. \\
$\mathit{nb\_credits\_anterieurs\_minimum}$ & $3$ & Arbitrage documenté explicitement dans
  \texttt{03-MODELE/\brk 04-cascade-\brk et-cold-start.md} : en dessous, le taux de remboursement
  du groupe repose sur trop peu d'observations et introduirait plus de bruit que de signal. \\
$\mathit{fraicheur\_maximale\_jours}$ & $7$ & Aucune source ni calibration documentée pour ce
  seuil précis -- à valider avec la coopérative. \\
\bottomrule
\end{tabular}

\subsection*{Logique}
\begin{formule}
\begin{verbatim}
SI les 4 conditions sont reunies
    ALORS mode = enrichi, modele = EBM_enrichi
SINON
    mode = socle_seul, modele = EBM_socle
    motif = raison precise de la bascule
\end{verbatim}
\end{formule}
Les quatre conditions (appartenance à un groupe, taille, historique de crédits soldés, fraîcheur
des features) sont vérifiées dans l'ordre de leur dépendance logique : on ne peut pas juger la
taille d'un groupe auquel on n'appartient pas, ni son historique avant d'avoir vérifié sa taille.

\subsection*{Pourquoi ce mécanisme}
45\,\% des dossiers du calibrage retenu sont des primo-emprunteurs : un système qui exigerait un
historique de groupe pour scorer serait inutilisable précisément là où le besoin est le plus fort.

\subsection*{Statut et rôle du modèle réel}
Cette règle choisit \emph{quel} modèle appeler (socle ou enrichi), elle ne calcule pas elle-même
de score. Les trois seuils restent des paramètres de politique ajustables par la coopérative,
indépendants de l'implémentation du modèle derrière chacun des deux ports.

\newpage

% ============================================================
\section{Montant recommandé par capacité de remboursement}
\label{sec:capacite}

\ecran{Résultat de scoring (montant recommandé) ; alerte agent si ce montant est réduit sur un
objet non divisible.}
\fichier{backend/\brk solida/\brk domain/\brk rules/\brk capacite\_remboursement.py::\brk{}montant\_maximal\_supportable}

\subsection*{Variables}

\begin{tabular}{@{}lL{9.5cm}@{}}
\toprule
Symbole & Signification \\
\midrule
$\mathit{revenu}$ & Revenu mensuel effectif (déclaré ou actualisé par l'agent) \\
$\mathit{duree}$, $\mathit{taux}$ & Durée et taux annuel du crédit demandé \\
\bottomrule
\end{tabular}

\subsection*{Constantes}

\begin{tabular}{@{}L{3.2cm}rL{7.4cm}@{}}
\toprule
Constante & Valeur & Pourquoi cette valeur \\
\midrule
$\mathit{ratio\_endettement\_maximal}$ & $0{,}33$ & \textbf{Sourcé, mais pas par FUCEC.} Étude de
  terrain CEF-MF Lomé 2025 (\texttt{PLAN 72H/\brk SOLIDA\_Complements\_et\_Strategie.md} §1.2) :
  ratio mensualité/revenu moyen de $23\,\%$ sur l'échantillon, rupture du risque d'impayé
  identifiée à $33\,\%$. Ni FUCEC-Togo (page publique des crédits, aucun ratio chiffré) ni le
  dispositif prudentiel BCEAO/UMOA (encadre le taux d'usure, pas un ratio d'endettement
  individuel) ne le documentent. C'est une étude d'un autre réseau togolais, pas encore confirmée
  par le comité de crédit de la coopérative elle-même. Voir
  \texttt{PLAN 72H/\brk DECISION\_MANQUANTE\_ratio\_endettement.md}. \\
\bottomrule
\end{tabular}

\subsection*{Formule}
\begin{formule}
\begin{align}
\mathit{mensualite\_unitaire} &= \mathit{mensualite\_actuarielle}(1,\ \mathit{duree},\ \mathit{taux}) \\
\mathit{montant\_maximal} &= \frac{\mathit{revenu} \times \mathit{ratio\_endettement\_maximal}}{\mathit{mensualite\_unitaire}} \\
\mathit{montant\_recommande} &= \min(\mathit{montant\_demande},\ \mathit{montant\_maximal})
\end{align}
\end{formule}
Réutilise la même fonction d'amortissement que le SOCLE et l'enrichi (linéarité en montant à
durée et taux fixés), plutôt que de dupliquer la formule financière.

\subsection*{Pourquoi cette formule}
Demande explicite d'un agent de crédit interrogé pendant la préparation terrain : recommander un
montant, mais avec une justification vérifiable variable par variable plutôt qu'un coefficient
sur le prêt précédent (mécanisme de crédit progressif, écarté par la décision terrain
« pas de coefficient fondé sur le prêt précédent »). Sans revenu connu, aucun plafond n'est
calculé -- décision terrain déjà actée, un revenu manquant ne réduit ni ne refuse rien
automatiquement.

\subsection*{Statut et rôle du modèle réel}
Indépendant du modèle de scoring : ce plafond ne dépend que de la demande et du revenu, pas de
la probabilité de défaut. Modifiable sans déploiement via \texttt{ratio\_endettement\_maximal}
dans la grille (\texttt{/parametrage/grille}) -- la coopérative peut l'ajuster dès qu'elle
dispose d'une vraie valeur.

\newpage

% ============================================================
\section{Tendance de l'épargne}
\label{sec:epargne}

\ecran{Dossier sociétaire (badge « en hausse » / « stable » / « en érosion » sur le graphique de
trajectoire d'épargne) ; alimente aussi la variable de scoring \texttt{tendance\_epargne\_12m}.}
\fichier{backend/\brk solida/\brk domain/\brk rules/\brk epargne.py::\brk{}tendance\_depuis\_croissance}

\subsection*{Constantes}

\begin{tabular}{@{}L{3.2cm}rL{6.9cm}@{}}
\toprule
Constante & Valeur & Pourquoi cette valeur \\
\midrule
$\mathit{SEUIL\_HAUSSE}$ & $10\,\%$ & \textbf{Aucune source.} Choisi à la main pour donner un
  libellé lisible à l'agent au guichet. \\
$\mathit{SEUIL\_EROSION}$ & $-5\,\%$ & \textbf{Aucune source.} Même choix, non symétrique par
  intuition (une érosion de l'épargne est jugée significative plus tôt qu'une hausse). \\
\bottomrule
\end{tabular}

C'est le seul cas de ce document où même l'ordre de grandeur est incertain -- contrairement au
taux BCEAO de la section~\ref{sec:echeance}, il n'existe pas de référence externe pour ce
découpage.

\subsection*{Formule}
\begin{formule}
\[
\mathit{tendance} =
\begin{cases}
\text{hausse} & \text{si } \mathit{croissance\_12m} > \mathit{SEUIL\_HAUSSE} \\
\text{érosion} & \text{si } \mathit{croissance\_12m} < \mathit{SEUIL\_EROSION} \\
\text{stable} & \text{sinon}
\end{cases}
\]
\end{formule}

\subsection*{Statut et rôle du modèle réel}
Cette fonction alimente à la fois l'affichage et une variable d'entrée du modèle. Un modèle réel
entraîné directement sur \texttt{croissance\_12m} en continu n'a pas besoin de ce découpage en 3
classes, qui perd de l'information par construction. À trancher avec le modèle réel : soit
recalibrer ces deux seuils sur des cohortes réelles, soit ne garder ce libellé que pour
l'affichage agent et passer la variable continue au modèle.

\newpage

% ============================================================
\section{Échéance mensuelle et taux d'endettement}
\label{sec:echeance}

\ecran{Formulaire « Nouvelle demande » (« Échéance mensuelle estimée ») ; alimente la variable de
scoring \texttt{ratio\_endettement}.}
\fichier{backend/\brk solida/\brk domain/\brk rules/\brk echeance.py::\brk{}calculer\_echeance\_mensuelle, calculer\_taux\_endettement}

\subsection*{Variables}

\begin{tabular}{@{}lL{9.5cm}@{}}
\toprule
Symbole & Signification \\
\midrule
$D$ & Montant demandé \\
$n$ & Durée en mois \\
\bottomrule
\end{tabular}

\subsection*{Constante}

\begin{tabular}{@{}L{3.2cm}rL{6.5cm}@{}}
\toprule
Constante & Valeur & Pourquoi cette valeur \\
\midrule
$i$ (taux mensuel) & $0{,}18/12$ (18\,\%/an) & Valeur de démonstration, nettement sous le plafond
  légal de 24\,\%/an (TAEG plafond BCEAO pour les SFD, UEMOA, depuis le 1\textsuperscript{er} juin
  2026), choisie en l'absence d'une table de taux réelle par produit de crédit. \\
\bottomrule
\end{tabular}

\subsection*{Formule (amortissement à annuité constante)}
\begin{formule}
\[
\mathit{échéance} = D \times \frac{i (1+i)^n}{(1+i)^n - 1}
\]
\end{formule}
\begin{formule}
\[
\mathit{taux\_endettement} = \frac{\mathit{charges\_mensuelles} + \mathit{échéance}}{\mathit{revenu\_mensuel}}
\]
\end{formule}

\subsection*{Pourquoi cette formule}
L'amortissement à annuité constante (intérêt dégressif sur le capital restant dû) est la méthode
que la réglementation BCEAO impose pour exprimer le TEG/TAEG des crédits dans l'UEMOA, et celle
qu'appliquent en pratique les coopératives d'épargne et de crédit togolaises -- pas un taux
forfaitaire sur le montant initial.

\subsection*{Statut et rôle du modèle réel}
Indépendant du modèle de scoring : le taux $i$ dépend d'une vraie table de taux par produit à
obtenir de la coopérative, pas d'un entraînement de modèle. Le taux d'endettement calculé ici
devient en revanche une variable d'entrée du modèle (section~\ref{sec:contributions-factices}),
dont le modèle réel apprendra le poids véritable.

\newpage

% ============================================================
\section{Politique de mot de passe}
\label{sec:motdepasse}

Incluse par exhaustivité : ce n'est pas un calcul de crédit, mais c'est une règle à seuils
numériques avec sa propre justification, comme tout ce qui précède.

\ecran{Connexion, changement de mot de passe.}
\fichier{backend/\brk solida/\brk domain/\brk rules/\brk mot\_de\_passe.py::\brk{}valider\_mot\_de\_passe}

\subsection*{Constantes}

\begin{tabular}{@{}L{3.2cm}rL{6.9cm}@{}}
\toprule
Constante & Valeur & Pourquoi cette valeur \\
\midrule
\texttt{LONGUEUR\_\brk{}MINIMALE} & 8 caractères & NIST SP 800-63B rev.\,4 : 8 caractères est le
  plancher \emph{SHALL} (obligatoire), quel que soit le nombre de facteurs d'authentification --
  15 n'est qu'une recommandation \emph{SHOULD} renforcée en l'absence de second facteur. Choix
  produit : rester praticable pour des agents de coopérative sans authentification
  multi-facteurs. \\
\texttt{LONGUEUR\_\brk{}MAXIMALE} & 64 caractères & Même norme : plafond \emph{SHALL} pour éviter les
  attaques par déni de service via des mots de passe démesurément longs. \\
\bottomrule
\end{tabular}

\subsection*{Règle (pas de formule)}
Aucune règle de composition (majuscule/chiffre/symbole) : la même norme l'interdit explicitement
-- ces règles poussent statistiquement vers des mots de passe prévisibles (substitutions
convenues, ex. « @ » pour « a »). Rejet uniquement sur la longueur et sur une liste de mots de
passe trop courants.

\subsection*{Statut et rôle du modèle réel}
Sans rapport avec le modèle de scoring.

\newpage

\section*{Récapitulatif}

\begin{longtable}{@{}L{4cm}L{3.4cm}L{3cm}L{3.2cm}@{}}
\toprule
Calcul & Fichier & Le modèle réel change... & ...mais pas \\
\midrule
\endhead
Score (PDO) & \texttt{scorecard.py} & $p$ \emph{et} les 3 paramètres d'échelle (recalibrage
  requis) & Les bornes d'affichage 300/850 \\
Invariant de cohérence & \texttt{scorecard.py} & Rien & La tolérance de 1 point \\
Décomposition en points & \texttt{scorecard.py} & Les contributions $f_j$ (deviennent réelles) &
  La formule elle-même \\
Contributions factices & \texttt{modele\_\allowbreak constant.py} & Remplacée entièrement & -- \\
Grille (tranche) & \texttt{grille.py} & $p$ & marge, lgd, multiplicateurs \\
Plafond progressif & \texttt{progressif.py} & $p$ & coefficient, plancher, modulation \\
Palier au prochain cycle & \texttt{progressif.py} & $p$ & idem + \texttt{nb\_cycles=1} \\
Conditions de réexamen & \texttt{progressif.py} & Rien & Toutes les cibles \\
Cascade (mode) & \texttt{cascade.py} & Rien (choisit quel modèle appeler) & Tous les seuils \\
Tendance d'épargne & \texttt{epargne.py} & À trancher (voir section~\ref{sec:epargne}) & -- \\
Échéance / endettement & \texttt{echeance.py} & Rien & Taux, indépendant du modèle \\
Mot de passe & \texttt{mot\_de\_\allowbreak passe.py} & Rien & Sans rapport avec le modèle \\
\bottomrule
\end{longtable}

\end{document}
